% =====================================================================
% 06-StabilityCriteria.tex
% Geometric Hallucination Taxonomy — Mathematical Spine
% Borealis S. Hedling — 2026
% =====================================================================

\section*{06. Unified Stability Criteria}

This module formalizes the unified instability framework used across the
Geometric Hallucination Taxonomy (GHT). Each operator—curvature, spectral,
holonomy, topology, and fusion—has its own instability criterion. Here we
combine them into a single stability structure that governs latent
trajectory behavior.

All definitions rely on the geometric foundations established in
00-Foundations.tex and the operator modules 01–05.

% ---------------------------------------------------------------------
\subsection*{6.1 Latent Trajectory}

Let $x(t)$ be a latent trajectory evolving on the manifold $M$:
\[
x : \mathbb{N} \to M.
\]

At each step $t$, the trajectory may exhibit geometric instability
according to one or more operator criteria.

% ---------------------------------------------------------------------
\subsection*{6.2 Operator Outputs}

For each operator, define the diagnostic output:

\begin{itemize}
\item Curvature: $K_t = K(\dot{x}(t), v(t))$,
\item Spectral: $\Delta_\lambda(t) = \lambda(t+1) - \lambda(t)$,
\item Holonomy: $H_{\gamma_t} = P_{\gamma_t} - I$,
\item Topology: $\tau_t = \tau(x(t))$,
\item Fusion: $F_t = F(x(t), y(t))$.
\end{itemize}

These quantities measure geometric behavior along the trajectory.

% ---------------------------------------------------------------------
\subsection*{6.3 Stability Thresholds}

Each operator has an associated threshold:

\[
\kappa_{\mathrm{thr}}, \quad
\sigma_{\mathrm{thr}}, \quad
h_{\mathrm{thr}}, \quad
\tau_{\mathrm{thr}} = n, \quad
f_{\mathrm{thr}}.
\]

A trajectory is stable with respect to an operator if its diagnostic
output remains below the corresponding threshold.

% ---------------------------------------------------------------------
\subsection*{6.4 Operator Stability Conditions}

\paragraph{Curvature Stability}
\[
|K_t| \le \kappa_{\mathrm{thr}}.
\]

\paragraph{Spectral Stability}
\[
|\Delta_\lambda(t)| \le \sigma_{\mathrm{thr}}.
\]

\paragraph{Holonomy Stability}
\[
\|H_{\gamma_t}\| \le h_{\mathrm{thr}}.
\]

\paragraph{Topological Stability}
\[
\tau_t = n.
\]

\paragraph{Fusion Stability}
\[
\|F_t\|_g \le f_{\mathrm{thr}}.
\]

% ---------------------------------------------------------------------
\subsection*{6.5 Unified Stability Region}

Define the unified stability region
\[
\Omega
= \Omega_{\kappa}
\cap \Omega_{\sigma}
\cap \Omega_{h}
\cap \Omega_{\tau}
\cap \Omega_{f}.
\]

A latent trajectory is globally stable if
\[
x(t) \in \Omega \quad \text{for all } t.
\]

% ---------------------------------------------------------------------
\subsection*{6.6 Unified Instability Condition}

A latent trajectory exhibits geometric instability if at least one
operator exceeds its threshold:

\[
\exists t \in \mathbb{N}, \quad
\Big(
|K_t| > \kappa_{\mathrm{thr}}
\;\lor\;
|\Delta_\lambda(t)| > \sigma_{\mathrm{thr}}
\;\lor\;
\|H_{\gamma_t}\| > h_{\mathrm{thr}}
\;\lor\;
\tau_t < n
\;\lor\;
\|F_t\|_g > f_{\mathrm{thr}}
\Big).
\]

This condition defines the onset of hallucination behavior.

% ---------------------------------------------------------------------
\subsection*{6.7 Persistent Instability}

Persistent instability occurs when instability recurs infinitely often:

\[
\{ t : x(t) \notin \Omega \} \text{ is infinite}.
\]

Persistent instability corresponds to sustained hallucination behavior.

% ---------------------------------------------------------------------
\subsection*{6.8 Operator Interaction}

Operators may interact. For example:

\begin{itemize}
\item Spectral cascades may induce curvature blowout,
\item Holonomy drift may distort fusion–field geometry,
\item Metric collapse may amplify spectral instability,
\item Fusion spikes may distort curvature directions.
\end{itemize}

Define the interaction map
\[
\mathcal{I} : \{K_t, \Delta_\lambda(t), H_{\gamma_t}, \tau_t, F_t\}
\to \{K_{t+1}, \Delta_\lambda(t+1), H_{\gamma_{t+1}}, \tau_{t+1}, F_{t+1}\}.
\]

Instability propagates through $\mathcal{I}$.

% ---------------------------------------------------------------------
\subsection*{6.9 Unified Hallucination Condition}

A latent trajectory exhibits hallucination behavior if
\[
\exists t_0 \in \mathbb{N}, \quad
x(t_0) \notin \Omega,
\]
and instability propagates through $\mathcal{I}$ to produce persistent
instability.

This condition unifies CDH, SCH, HDH, TIH, and FFH.

% ---------------------------------------------------------------------
\subsection*{6.10 Discrete Approximation}

In numerical settings, unified stability is checked by evaluating each
operator’s discrete approximation at each step and verifying whether
$x(t)$ remains in $\Omega$.

The Python toolkit implements unified stability checks in
\texttt{/tests/test\_stability.py}.

% =====================================================================
% End of 06-StabilityCriteria.tex
% =====================================================================
